% ─────────────────────────────────────────────────────────────────────────────
\begingroup
\setlength{\intextsep}{8pt plus 2pt minus 2pt}
\setlength{\textfloatsep}{10pt plus 2pt minus 4pt}
\setlength{\floatsep}{10pt plus 2pt minus 4pt}
\captionsetup[figure]{skip=5pt}
\newcommand{\sdsfig}[1]{%
  \includegraphics[width=\linewidth,height=0.69\textheight,keepaspectratio]{#1}%
}

This chapter provides the critical design artifacts for the BookForMe project. It details the system's design decisions, hybrid architectural patterns, data structures, and the workflow logic required to build a robust service booking application.

% ─────────────────────────────────────────────────────────────────────────────
\section{Introduction}
\label{sec:sds-intro}

\subsection{Purpose}
The purpose of this Software Design Specification (SDS) is to provide a comprehensive architectural overview of the \textbf{BookForMe} platform. It establishes a single source of truth for all technical decisions, detailing the accurate architecture from the database schema to the conversational AI pipelines.

\subsection{Scope}
The BookForMe system connects service providers (Vendors) and consumers (Users) in Karachi, Pakistan through a centralized platform. It solves the inefficiency of manual booking by automating the entire booking and payment verification process. The system includes:
\begin{itemize}
    \item \textbf{A User Mobile App:} A React Native (Expo) interface for users to browse, book, and engage socially.
    \item \textbf{An Automated WhatsApp Agent:} An intelligent assistant that handles direct bookings conversations on behalf of the vendor using natural language in English or Roman Urdu.
    \item \textbf{A Vendor Dashboard:} A React web interface for business owners to view their schedule, manage slots, and give final approval on bookings.
    \item \textbf{A Centralized Backend:} A Python FastAPI layer that keeps data consistent, ensuring that a slot booked via WhatsApp is instantly updated on the App.
\end{itemize}

\subsection{Definitions and Acronyms}
\begin{itemize}
    \item \textbf{Neuro-Symbolic Agent Pipeline:} An architecture utilizing LangGraph and Qwen 3 32B to parse natural language into strict intents, while deterministic code handles database writes to prevent hallucinated bookings.
    \item \textbf{Rich Client (Mobile App):} A fully-featured application installed on the phone.
    \item \textbf{Conversational Client (WhatsApp):} A text-based Webhook interface connected to an AI Receptionist.
    \item \textbf{OCC (Optimistic Concurrency Control):} A strategy using Firestore transactional locks to prevent double-booking across simultaneous users.
    \item \textbf{OCR (Optical Character Recognition):} Technology utilizing Gemini Vision to automatically extract and verify amounts from payment screenshots.
\end{itemize}

% ─────────────────────────────────────────────────────────────────────────────
\section{System Architecture}
\label{sec:sds-arch}

\subsection{Architectural Design}
The system follows a \textbf{Service-Oriented Architecture (SOA)} with a clear separation between the Presentation, Business Logic, and Persistence layers. Both the Rich Client and the Conversational Client converge on a single backend and Google Cloud Firestore database, ensuring a "Single Source of Truth." Dependencies are strictly unidirectional: Clients depend on the Backend, and the Backend depends on the Database.

\subsection{System Architecture Diagram}
The diagram below illustrates the high-level data flow across the system's three tiers.
\begin{itemize}
    \item \textbf{Synchronous API Flow:} The Mobile Application and Vendor Dashboard communicate directly with the Central Backend API via JSON/HTTPS requests.
    \item \textbf{Asynchronous Event Flow:} The WhatsApp API operates via webhooks, triggering an asynchronous event routed to the LangGraph AI Receptionist.
\end{itemize}

\begin{figure}[H]
    \centering
    \IfFileExists{images/System Architecture.png}{
        \sdsfig{images/System Architecture.png}
    }{\fbox{System Architecture Diagram Missing}}
    \caption{High-Level System Architecture}
    \label{fig:sds-system-arch}
\end{figure}

% ─────────────────────────────────────────────────────────────────────────────
\section{Data Model (Data Design)}
\label{sec:sds-data}

\subsection{Database Strategy}
Data persistence is handled by \textbf{Google Cloud Firestore}. We employ a "Reference-Based" schema strategy to balance the flexibility of NoSQL with the need for structured relationships. The design strictly separates the transactional domain (vendors, slots, payments) from social or profile elements to ensure data integrity.

\subsection{Entity Descriptions}
Table~\ref{tab:entities} describes the primary data collections deployed in Firestore.

\begin{table}[H]
    \centering
    \caption{Domain Entity Descriptions}
    \label{tab:entities}
    \small
    \begin{tabular}{|p{3cm}|p{5cm}|p{6cm}|}
    \hline
    \textbf{Collection} & \textbf{Key Attributes} & \textbf{Description \& Purpose} \\ \hline
    \textbf{users} & phone, role, stats & Maintains core identity, authentication references, and gamification profiles. \\ \hline
    \textbf{vendors} & operating\_hours, revenue & Represents the business entity and operational constraints. \\ \hline
    \textbf{services} \& \textbf{resources} & pricing, capacity & Sellable units (e.g., Padel) mapped to specific real-world physical assets (e.g., Court 1). \\ \hline
    \textbf{slots} & status, hold\_expires\_at & The core inventory ledger representing blockable time intervals. \\ \hline
    \textbf{payments} & screenshot\_url, ocr\_verified\_amount & Audit trail of transactions, separated to ensure OCR verification operates safely before bookings are confirmed. \\ \hline
    \end{tabular}
\end{table}

\subsection{NoSQL Schema Diagram}
The Entity Relationship Diagram delineates references between entities, mapping out the foreign keys (e.g., \texttt{vendor\_id}) that correlate disparate collections.

\begin{figure}[H]
    \centering
    \IfFileExists{images/ERD.png}{
        \sdsfig{images/ERD.png}
    }{\fbox{ERD Diagram Missing}}
    \caption{Entity Relationship Diagram defining referential structures.}
    \label{fig:sds-erd}
\end{figure}

% ─────────────────────────────────────────────────────────────────────────────
\section{Agentic Booking Process Pipeline}
\label{sec:sds-agent-pipeline}

A core functionality is the AI-driven booking process managed by a LangGraph temporal state machine. It manages stateful conversations, enforcing a strict validation pipeline.

\subsection{Pipeline Logic}
\begin{enumerate}
    \item \textbf{Ingest \& Negotiation:} The Agent loops with the user to clarify intents (e.g., \texttt{GREETING}, \texttt{TRANSACTION}) and specific entities (date, time, service).
    \item \textbf{Transaction (Locking):} Using OCC, the system creates a \texttt{HOLD} booking state with a 10-minute expiry window.
    \item \textbf{Intelligent Gatekeeper:} Gemini Vision OCR extracts the payment screenshot amount. If invalid, the Agent auto-rejects quietly, asking for a re-upload.
    \item \textbf{Human-in-the-Loop:} Valid receipts reach the Vendor Dashboard for a final localized authorization.
\end{enumerate}

\begin{figure}[H]
    \centering
    \IfFileExists{images/FYP - Seq Dig.png}{
        \sdsfig{images/FYP - Seq Dig.png}
    }{\fbox{Sequence Diagram Missing}}
    \caption{AI Booking \& Payment Pipeline}
    \label{fig:sds-seq-dig}
\end{figure}

% ─────────────────────────────────────────────────────────────────────────────
\section{Technical Details \& Workflows}
\label{sec:sds-technical-workflows}

\subsection{Key Algorithms}
\subsubsection{Algorithm 1: Intent Parsing \& Action Routing}
\textbf{Purpose:} To translate unstructured natural language into deterministic database actions without hallucination.
\begin{itemize}
    \item \textbf{Input:} Raw WhatsApp Message string + User Session State.
    \item \textbf{Logic:} Qwen 3 32B classifies intent. The state machine "routes" this to a deterministic Python node ensuring the LLM does not generate database queries directly.
    \item \textbf{Output:} Validated JSON Payload or Clarification Prompt.
\end{itemize}

\subsubsection{Algorithm 2: Optimistic Concurrency Control (OCC)}
\textbf{Purpose:} To prevent race conditions (double bookings) in a distributed environment.
\begin{itemize}
    \item \textbf{Input:} \texttt{service\_id}, \texttt{time\_slot}, \texttt{user\_id}.
    \item \textbf{Logic:}
    \begin{enumerate}
        \item \textbf{Read:} Fetch the current slot status.
        \item \textbf{Verify:} Ensure status is \texttt{AVAILABLE}.
        \item \textbf{Write:} Atomically update to \texttt{LOCKED} with \texttt{hold\_expires\_at}. Fails entirely if modified by another request.
    \end{enumerate}
    \item \textbf{Output:} Transaction Success or Conflict Error (triggers alternative slot suggestions).
\end{itemize}

\subsection{Workflow Visualizations}
\FloatBarrier

\Needspace{0.82\textheight}
\subsubsection{Workflow 1: The Agentic Booking Experience}
This covers the WhatsApp conversational flow, focusing heavily on the Error Correction Loop wherein rejected OCR receipts autonomously trigger a re-prompt.
\begin{figure}[H]
    \centering
    \IfFileExists{images/Activity Diagram- Whatsapp Booking flow.png}{
        \sdsfig{images/Activity Diagram- Whatsapp Booking flow.png}
    }{\fbox{Activity Diagram: Agent Flow Missing}}
    \caption{Activity Diagram: Agent/WhatsApp Booking Flow}
\end{figure}

\FloatBarrier
\Needspace{0.82\textheight}
\subsubsection{Workflow 2: The App User Experience}
This underscores the React Native client's Dual-Search Capability, permitting classic visual filtering alongside natural language queries pointing to the identical underlying slot ledger.
\begin{figure}[H]
    \centering
    \IfFileExists{images/AD- App flow booking process.png}{
        \sdsfig{images/AD- App flow booking process.png}
    }{\fbox{Activity Diagram: App Flow Missing}}
    \caption{Activity Diagram: Mobile App Search \& Book}
\end{figure}

\FloatBarrier
\Needspace{0.82\textheight}
\subsubsection{Workflow 3: The Vendor Management Experience}
This maps the Vendor integration, indicating that rejections trigger a sweeping cascade releasing \texttt{LOCKED} slots back to \texttt{AVAILABLE}.
\begin{figure}[H]
    \centering
    \IfFileExists{images/FYP - Vend Activity Dig.png}{
        \sdsfig{images/FYP - Vend Activity Dig.png}
    }{\fbox{Activity Diagram: Vendor Flow Missing}}
    \caption{Activity Diagram: Vendor Dashboard \& Approvals}
\end{figure}

\endgroup

% ─────────────────────────────────────────────────────────────────────────────
\section{Security Implementation}
\label{sec:sds-security}
\begin{itemize}
    \item \textbf{Authentication:} Managed via Firebase Auth, using the standard JSON Web Tokens mapped against user claims.
    \item \textbf{Role-Based Access Control (RBAC):} FastAPI dependencies validate claims (Admin vs. Vendor vs. Customer) to prevent unauthorized API calls.
    \item \textbf{AI Guardrails:} The deterministic LangGraph nodes create an inherent prompt injection firewall. Because the LLM lacks arbitrary read/write permissions directly reaching Firestore, unauthorized transactions are logically impossible.
    \item \textbf{Data Privacy:} Payment screenshots populate isolated Firebase Storage buckets. They are served entirely via ephemeral, time-limited signed URLs only readable by authorized vendor proxies.
\end{itemize}

%%% Local Variables:
%%% mode: latex
%%% TeX-master: "../report"
%%% End:
